\section{Introduction}
\subsection{Objectif}
Nous étudions l'apport de la contrainte d'ordre pour identifier des régimes
homogènes dans une séquence. Fisher est comparé aux K-means et à la CAH de Ward
sur \file{sequencesimu}, puis sur \file{aiguillage} à quatre classes.
Une classe doit ici former un intervalle : des groupes proches dans l'espace
peuvent correspondre à des régimes temporels différents.

\subsection{Données et démarche}
\begin{table}[H]
\centering\small
\begin{tabular}{lrrl}
\toprule
Fichier & Observations $n$ & Variables $p$ & Référence disponible \\
\midrule
\file{sequencesimu.txt} & 160 & 2 & Aucune étiquette fournie \\
\file{aiguillage.txt} & 140 & 552 & Colonne 553 : quatre classes \\
\bottomrule
\end{tabular}
\caption[Données analysées]{Données analysées. Les lignes restent dans leur ordre d'origine.}
\label{tab:donnees}
\end{table}

La colonne de référence de \file{aiguillage} est séparée des variables avant
tout ajustement. Les 552 coordonnées contribuent au calcul des distances.

La démarche comporte trois étapes : définir et vérifier le critère de Fisher,
choisir le nombre de segments sur les données simulées, puis évaluer les
partitions réelles. Les expériences d'initialisation et de changement d'échelle
évaluent la sensibilité des partitions.
